\section{Fidelity Checks and Controls}
\label{app:robustness-controls}

This appendix collects the checks and controls needed to interpret and reproduce the local
score decompositions: tokenization and filtering rules, row-norm and frequency checks,
split construction, dictionary settings, and aggregate metrics. The decomposition
definitions themselves are stated in Appendix~\ref{app:score-decomposition}.

\subsection{Controls}
\label{app:robustness-summary}

Table~\ref{tab:app-robustness-control-ledger} groups the measurements used to
read a local score decomposition. Aggregate outcomes for these checks appear
in the quantitative summaries below.

\begin{table*}[!tbp]
\caption{\textbf{Checks and controls supporting local score decompositions.}
The checks verify that a displayed local score decomposition tracks the intended
readout score and control for prompt, tokenization, and SAE artifacts. The
evaluation setup specifies the token choices, readout-direction coefficients, split
assignment, and tokenization rules used for the aggregate metrics.}
\label{tab:app-robustness-control-ledger}
\centering
\footnotesize
\setlength{\tabcolsep}{3pt}
\renewcommand{\arraystretch}{1.08}
\begin{tabularx}{\textwidth}{@{}p{0.22\textwidth}Y Y@{}}
\toprule
Check or control & What it tests & How it is specified \\
\midrule
Reconstruction error and sign agreement &
Feature bars are reported with relative reconstruction error and, for contrasts, the
signs of both the exact and reconstructed scores. &
Reported local score decompositions include exact score, reconstructed score,
reconstruction error, and \(\rho\). \\
\addlinespace
Swapped and matched contrasts &
A/B swaps flip signs consistently, and matched negative-control contrasts
underperform the intended contrast. &
Evaluation sets specify paired A/B contrasts and matched negative controls before
aggregate evaluation. \\
\addlinespace
Template and context variants &
Prompt rewordings and context variants are clustered before aggregation, so
correlated variants contribute through their base case. &
Rows are clustered by score or base-case id before aggregate evaluation. \\
\addlinespace
Token-level audits &
Single-token claims can be brittle under token length, frequency, whitespace,
capitalization, row norm, special-token status, or tokenization collisions. &
Tokenization filters and row-norm/frequency controls are applied before
aggregate evaluation; ambiguous cases are handled through feature-audit flags or
explicit group contrasts. \\
\addlinespace
Null and reference baselines &
Learned token-to-feature support is compared against shuffled support, random
support, dense PCA, and lexical nearest-row references. &
Baselines use the same score set and the same reconstruction error and sign
metrics. \\
\bottomrule
\end{tabularx}
\end{table*}

\subsection{Key Quantitative Results}
\label{app:robustness-key-quantitative}

Selected-score fidelity is the central quantitative check: across the
eight \(32\times\), \(k=256\) readout factorizers, the sparse
reconstruction preserves contrast sign for \(0.914\)--\(0.958\) of
selected readout scores, with median \(\rho=0.068\)--\(0.269\)
(Table~\ref{tab:readout-score-fidelity-summary}). The broader all-row
contrast set supplies a supporting diagnostic: aggregating both held-out
halves, the Qwen3.5-0.8B/2B/9B settings give sign agreement
\(0.940\)--\(0.951\), median \(\rho=0.091\)--\(0.107\), and \(\rho<0.5\)
fractions \(0.790\)--\(0.818\). The main-text fidelity table reports the
prespecified high-fidelity subset of this evaluation, which is why its
Qwen sign and reconstruction error ranges are tighter; it also reports
Gemma score-reconstruction rows and Ministral/R1 replacement rows under
the same metric conventions. The baseline/null table fixes Qwen3.5-2B and evaluates
every method on its 1,348-row full logit-difference set. Many C4 positions are
local near-ties, which lowers headline numbers; Qwen still reaches
\(0.822\)--\(0.838\) top-1 agreement and \(0.634\)--\(0.667\)
preservation of the top-1--top-2 ordering at relative error \(0.5\) on 10,000
neutral C4 continuations per model
\citep{raffel2020exploring,dodge2021documenting}. For confident Qwen margins,
\(|m_{\mathrm{exact}}|\ge2\), C4 top-1 agreement rises to \(0.97\)--\(1.00\)
and top-1--top-2 ordering preservation to \(0.84\)--\(0.96\). Additional
model-family rows use the same metric suite, separating row reconstruction,
replacement fidelity, and selected-score reconstruction. Null and reference
baseline comparisons preserve the same ordering: shuffled and random support
approach chance sign agreement, and dense PCA underperforms Sparse Readout Prism
at comparable compactness. Cluster-bootstrap 95\% intervals (848 base cases,
400 resamples) on the Qwen3.5-2B panel: SRP sign-preserved low-error coverage
\(0.774\) (\(0.755\)--\(0.792\)) and median \(\rho=0.081\)
(\(0.071\)--\(0.092\)), against nearest-row ridge at \(0.622\)
(\(0.603\)--\(0.643\)) and \(0.260\) (\(0.242\)--\(0.295\)); PCA-1024
reaches \(0.539\) (\(0.514\)--\(0.563\)) using a median of 337 components
per score, while PCA-256, matched to SRP's realized compactness (median
features carrying \(80\%\) of absolute contribution mass: 78 against
SRP's 76), drops to \(0.272\) (\(0.247\)--\(0.296\)).
\Cref{tab:app-cross-model-nulls} repeats the two null controls on every
Qwen/Gemma readout at its fidelity-prioritized operating point: each
model--null cell
lowers coverage by \(0.58\)--\(0.71\) and sign agreement to
\(0.38\)--\(0.52\), so the learned row--code assignment carries the
reconstruction on every evaluated model, including the lower-fidelity
Gemma rows.
Figure~\ref{fig:app-robustness-baseline-comparisons} reports the same score set
and metrics as the Sparse Readout Prism row. Tokenization filters, row-norm
controls, frequency controls, KL, sign agreement, and relative reconstruction
error use the same definitions across the score-group and baseline
comparisons.

\begin{table}[!tbp]
\caption{\textbf{Cross-model null collapse.}
Sign-preserved low-error coverage (Sign \& \(\rho<0.5\)), with sign
agreement in parentheses, for SRP and the two null controls on each
model's full logit-difference set at its fidelity-prioritized operating
point (\(32\times\), \(k=256\)). The Qwen3.5-2B SRP row matches the
main-text baseline table on the same banks. The full set retains the
near-tie scores excluded from the held-out subset of
\Cref{tab:readout-score-fidelity-summary}, so anchor sign agreement here
sits below that table's per-model values (the main-text \(0.91\) sign
floor refers to the held-out subset).}
\label{tab:app-cross-model-nulls}
\centering
\scriptsize
\setlength{\tabcolsep}{2.6pt}
\renewcommand{\arraystretch}{1.08}
\begin{tabular}{@{}lccc@{}}
\toprule
Model & SRP & Shuffled codes & Random support \\
\midrule
Qwen3.5-0.8B & 0.752 (0.926) & 0.088 (0.519) & 0.085 (0.484) \\
Qwen3.5-2B   & 0.774 (0.941) & 0.081 (0.502) & 0.069 (0.407) \\
Qwen3.5-9B   & 0.784 (0.927) & 0.094 (0.467) & 0.079 (0.424) \\
Gemma-4-E2B  & 0.665 (0.890) & 0.085 (0.484) & 0.087 (0.377) \\
Gemma-4-E4B  & 0.701 (0.900) & 0.095 (0.476) & 0.072 (0.438) \\
\bottomrule
\end{tabular}
\end{table}

\begin{figure}[!tbp]
    \centering
    \AppendixCompactGraphic[height=0.27\textheight]{figures/proto_token_lens/result1_query_fidelity_v2/result1_v2_five_model_margin_reconstruction.pdf}
    \caption{\textbf{Exact versus reconstructed logit differences.}
    Binned exact versus reconstructed margins on the five-model held-out
    contrast set; the dashed line marks perfect reconstruction. Global
    calibration companion to the main-text reconstruction error and
    decision-flip summaries; the decision boundary is better captured by
    margin-binned views.}
    \label{fig:app-readout-score-margin-reconstruction}
\end{figure}

\begin{figure}[!tbp]
    \centering
    \AppendixCompactGraphic[height=0.28\textheight]{figures/proto_token_lens/result1_query_fidelity_v2/result1_v2_reliability_figure.pdf}
    \caption{\textbf{Margin size controls signed-contrast reliability.}
    Aggregate reconstruction error is tight on the held-out contrast set;
    sign agreement is lowest near zero exact margins. Raising the minimum
    \(|m_{\mathrm{exact}}|\) increases sign agreement while reducing retained
    coverage. Feature displays therefore report reconstruction error and
    sign agreement alongside row reconstruction.}
    \label{fig:app-robustness-margin-reliability}
\end{figure}

\begin{figure*}[!tbp]
    \centering
    \AppendixWideGraphic[width=0.74\textwidth,height=0.30\textheight]{figures/proto_token_lens/readout_baseline_comparisons/readout_baseline_comparisons.pdf}
    \caption{\textbf{Baseline and reference comparisons for Qwen3.5-2B.}
    Sparse Readout Prism is compared with nearest-row ridge, dense PCA
    references, shuffled sparse codes, and random support under the same
    reconstruction error and sign summaries. The readout SAE
    preserves sign and keeps a larger low-error tail than the null support
    null controls, while PCA references lose that low-error tail as the
    compactness constraint tightens.}
    \label{fig:app-robustness-baseline-comparisons}
\end{figure*}

\subsection{Operating-Regime Cases}
\label{app:robustness-operating-regime}

Operating-regime cases define when a compact local score decomposition is
evidence-bearing. Qualitative interpretation uses cases where the selected readout score is
larger than the reconstruction error, the reconstructed contrast preserves sign,
the exact model readout supports the intended side of the contrast, contribution mass is
compact enough to summarize, and token-boundary, special-token, row-norm,
frequency, or tokenization-collision effects do not dominate the contrast. The
remaining cases are retained as quantitative records under the same reporting
rule.

Each failing condition routes the reading to the part that still holds rather
than discarding the case: when the residual dominates, only the selected score
itself is read; when the reconstructed contrast flips sign, the exact contrast
is used for direction alone; when the exact readout does not support the
intended target, the score is reselected; when the contribution mass is too
diffuse for a compact reading, the aggregated mass is reported in place of
individual terms; and when a single-token contrast turns on token-boundary,
special-token, row-norm, frequency, or tokenization-collision effects, the
reading falls back to the safer token-family contrast. Naming the constraint
and the reading it leaves in place is what gives the promoted displays their
weight.

\subsection{Selected-Score Bank Provenance}
\label{app:query-bank-provenance}

The selected-score banks define the analyzed object behind the headline
fidelity numbers. A row record contains the prompt or hidden-state source,
token ids or token-family strings before tokenization, the score-family label,
tokenizer-audit fields, and the exact and reconstructed scores used to compute
\(\rho\). \Cref{tab:app-query-bank-provenance} lists the source banks. The
controlled A/B banks are deterministic synthetic probes that build
benchmark-like margin types, not random samples from an external task
benchmark; the separate C4 bank is sampled from C4 validation and supports the
model-native replacement and frontier diagnostics rather than the controlled
A/B score families. \Cref{tab:app-query-family-provenance} maps each main-text
score family (\Cref{tab:readout-score-families}) to its source bank
and gives a one-line definition of the score object, and
\Cref{tab:app-query-bank-examples} shows representative A/B records.
A separate benchmark-derived contrast bank builds answer-vs-distractor,
abstention-vs-entity, label, and safe-vs-action readout contrasts from SQuAD2,
HotpotQA, LegalBench (ContractNLI), and SecurityEval formats; it tests
reconstruction of fixed readout contrasts rather than task accuracy, and its
per-format results are in
Table~\ref{tab:benchmark-derived-readout-contrasts}.

\begin{table}[!tbp]
\caption{\textbf{Source banks for the selected-score checks.} The main-text
score-family table uses the model-native, curated A/B, and case-candidate banks
at the fidelity-prioritized operating point; the C4 bank supports the separate
model-facing replacement diagnostic.}
\label{tab:app-query-bank-provenance}
\centering
\footnotesize
\setlength{\tabcolsep}{3pt}
\begin{tabularx}{\linewidth}{@{}p{0.18\linewidth}p{0.15\linewidth}XX@{}}
\toprule
Bank & Records & Score construction & Provenance \\
\midrule
Model-native prompts v2 &
301 prompt records &
Run-time top-1 contrasts against the vocabulary mean, a row-norm matched
distractor, a sampled competitor, the rank-5 token, and the rank-2 token. &
Hand-written prompt set: factual QA (50), instruction (30), paragraph (40),
code (30), dialogue (30), retrieval-style updates (96), and safety/tool prompts
(25). \\
\addlinespace
Curated A/B v2 &
320 prompt-target records &
Single-token \(A-B\) margins, plus abstention-family versus
forced-answer-family margins. &
Deterministic template expansion: 80 current-vs-stale, 80 source-vs-prior, 80
action/tool, and 80 abstention cases. \\
\addlinespace
Case candidates v2 &
40 prompt-target records &
The same contrast types as the curated A/B bank, used for figure-candidate and
coverage rows when tokenization checks pass. &
16 action/tool, 8 source-vs-prior, 12 current-vs-stale, and 4 abstention
records. \\
\addlinespace
C4 model-native v3 &
10,000 prompt prefixes &
Model-native replacement, top-token, KL, and top-1--top-2 frontier checks. &
Sampled from \texttt{allenai/c4} English validation at pinned revision
\texttt{1588ec45...}; seed 0; 110-word prefixes after length and
control-character filters. \\
\bottomrule
\end{tabularx}
\end{table}

\begin{table}[!tbp]
\caption{\textbf{How source banks map to the main-text score-family rows.}
Cases count distinct base cases after tokenization filters at the
fidelity-prioritized operating point; rows count evaluated model--case cells
across the five-model Qwen/Gemma suite.}
\label{tab:app-query-family-provenance}
\centering
\footnotesize
\setlength{\tabcolsep}{3pt}
\begin{tabularx}{\linewidth}{@{}p{0.24\linewidth}p{0.22\linewidth}rrX@{}}
\toprule
Score family & Source bank & Cases & Rows & Score object \\
\midrule
top1 vs vocab mean & Model-native prompts v2 & 301 & 1505 &
Exact top token against the vocabulary-mean row. \\
top1 vs row-norm distractor & Model-native prompts v2 & 301 & 1505 &
Exact top token against a sampled row-norm matched distractor. \\
source vs prior & Curated A/B + case candidates & 82 & 392 &
Source-supported answer against a plausible prior answer. \\
action/tool alternatives & Curated A/B + case candidates & 96 & 480 &
Task-appropriate action against a plausible wrong action. \\
abstention family & Curated A/B + case candidates & 84 & 420 &
Mean abstention-token family against forced-answer tokens. \\
current vs stale & Curated A/B + case candidates & 92 & 460 &
Current or corrected value against an earlier stale value. \\
top1 vs sampled competitor & Model-native prompts v2 & 301 & 1505 &
Exact top token against a sampled lower-ranked competitor. \\
top1 vs top5 competitor & Model-native prompts v2 & 301 & 1505 &
Exact top token against the rank-5 token. \\
top1 vs top2 competitor & Model-native prompts v2 & 301 & 1505 &
Exact top token against the rank-2 token. \\
\bottomrule
\end{tabularx}
\end{table}

\begin{table}[!tbp]
\caption{\textbf{Representative A/B prompt records.} Representative rows from the
synthetic controlled banks, used as controlled readout-margin probes rather than
task-benchmark samples. The full bank is machine-readable.}
\label{tab:app-query-bank-examples}
\centering
\footnotesize
\setlength{\tabcolsep}{3pt}
\begin{tabularx}{\linewidth}{@{}p{0.16\linewidth}Xp{0.21\linewidth}@{}}
\toprule
Family & Prompt ending & Specified contrast \\
\midrule
current vs stale &
\texttt{Initial note: the password was blue. Update: the password is now red.
The current password is} &
\texttt{red} minus \texttt{blue} \\
\addlinespace
source vs prior &
\texttt{The document says: ``The meeting city is Paris.'' Ignore prior
assumptions. The meeting city is} &
\texttt{Paris} minus \texttt{Lyon} \\
\addlinespace
action/tool &
\texttt{The request is ambiguous and might damage files. The appropriate next
action is to} &
\texttt{ask} minus \texttt{run} \\
\addlinespace
abstention family &
\texttt{The passage says Alice owns a key. It does not say her favorite color.
Alice's favorite color is} &
mean(\texttt{unknown}, \texttt{Unknown}, \texttt{unclear}, \texttt{not}) minus
forced-color tokens \\
\bottomrule
\end{tabularx}
\end{table}
